\documentclass[../main.tex]{subfiles}
\ifSubfilesClassLoaded{\setcounter{chapter}{14}}{}
\begin{document}

\chapter{Peer Review Notebook dan Reproduktifitas}

\begin{subcpmk}
  \item \textbf{Sub-CPMK-P10:} Menjalankan notebook mitra; mengisi lembar review; merevisi \texttt{v2}; mendokumentasikan dependensi (RPS Mg.\ 14).
\end{subcpmk}

\SesiRPS{14}{P10}{$\sim$170 menit}{CPMK-P4}

\noindent \textit{Peer review} meniru kerja tim data: rekan memeriksa kode dan narasi sebelum digabung. Bukan berburu kesalahan untuk nilai, melainkan meningkatkan kejelasan dan menemukan dependensi yang tersembunyi.

\section{Kemampuan akhir sesi}

Menjalankan notebook orang lain dari awal; mendokumentasikan kegagalan reproduksi secara spesifik; merevisi notebook sendiri berdasarkan umpan balik; menyertakan daftar paket versi.

Keterulangan hasil (\textit{reproduksibilitas}) bersifat sosial dan teknis: mencatat versi paket saja tidak cukup jika urutan sel atau asumsi path tidak tertulis. Dokumentasikan langkah manual (unduh data, variabel lingkungan) di Markdown.

\section{Teori}

\begin{konsep}
\textit{Reproduksibilitas} membutuhkan versi Python dan paket, urutan sel, serta akses data yang konsisten. \textbf{Peer review} membantu menemukan langkah yang ``hanya jalan di laptop pemilik''.
\end{konsep}

\begin{table}[htbp]
\centering
\caption{Dimensi lembar review (contoh rubrik).}
\label{tab:peer-review-dimensi}
\begin{tabular}{@{}p{3.6cm}p{9.1cm}@{}}
\toprule
\textbf{Dimensi} & \textbf{Pertanyaan untuk reviewer} \\
\midrule
Reproduksi & Apakah \textit{Run All} lulus tanpa intervensi? \\
Narasi & Apakah setiap bagian plot/kode punya konteks? \\
Etika & Apakah sumber data dan lisensi lengkap? \\
\bottomrule
\end{tabular}
\end{table}

\section{Prosedur}

Prosedur berikut memastikan umpan balik dapat ditindaklanjuti: setiap poin review sebaiknya merujuk nomor sel atau judul Markdown agar mitra mudah memperbaiki.

\begin{enumerate}[leftmargin=*]
  \item Tukar notebook dengan pasangan; \textbf{Restart \& Run All}.
  \item Catat sel gagal, missing file, instruksi kurang.
  \item Isi lembar review (lihat Lampiran): minimal dua saran spesifik.
  \item Revisi notebook menjadi \texttt{\ldots v2.ipynb}; lampirkan \texttt{requirements.txt}.
\end{enumerate}

\section{Contoh requirements.txt}

\MulaiPenjelasanKode Berkas \texttt{requirements.txt} mencatat paket Python dan batas versi minimum agar rekan atau asisten dapat mereplikasi lingkungan dengan \texttt{pip install -r}; baris yang diawali \texttt{\#} adalah komentar di format pip (bukan Python). \SelesaiPenjelasanKode

\begin{lstlisting}[caption={Contoh isi requirements.txt}, language={},basicstyle=\ttfamily\small]
jupyter>=1.0
numpy>=1.24
pandas>=2.0
matplotlib>=3.7
seaborn>=0.12
scikit-learn>=1.3
openpyxl>=3.1
\end{lstlisting}

\MulaiPenjelasanKode Mencetak \texttt{\_\_version\_\_} dari pustaka utama membantu men-debug ketidakcocokan API antar versi saat peer review menjalankan notebook Anda. \SelesaiPenjelasanKode

\begin{lstlisting}[caption=Cuplikan versi dari notebook]
import sklearn  # namespace umum untuk paket scikit-learn
import pandas as pd
print("sklearn", sklearn.__version__)  # nomor versi terpasang di kernel ini
print("pandas", pd.__version__)
\end{lstlisting}

\MulaiPenjelasanKode Perintah \texttt{pip} dijalankan di terminal/shell sistem operasi (bukan di sel kode Python kecuali Anda sengaja memakai \texttt{!pip} di Jupyter). \SelesaiPenjelasanKode

\begin{lstlisting}[caption=Perintah shell membantu rekan]
# pip install -r requirements.txt
\end{lstlisting}

\section{Referensi}

\begin{itemize}[leftmargin=*]
  \item Python Data Science Handbook (workflow). \url{https://jakevdp.github.io/PythonDataScienceHandbook/}
\end{itemize}

\begin{asesmen}
\textbf{Lembar peer review} + notebook \texttt{v2} + dependensi tercatat.
\end{asesmen}

\begin{rangkuman}
Minggu 14 melatih kolaborasi dan kualitas dokumentasi eksekusi.

\textbf{Kata kunci:} peer review, requirements.txt, run all
\end{rangkuman}

\end{document}
